\newcommand{\HeadlineRows}{3 & 3 & 1 & 0.059 & 0.033 & 0.030 & 1.80$\times$\\
6 & 3 & 1 & 0.409 & 0.155 & 0.157 & 2.64$\times$\\
12 & 3 & 1 & 0.730 & 0.240 & 0.266 & 3.05$\times$\\
12 & 3 & 2 & 1.403 & 0.316 & 0.312 & 4.43$\times$\\
12 & 12 & 1 & 2.687 & 0.804 & 1.090 & 3.34$\times$\\
12 & 48 & 1 & 11.177 & 3.188 & 6.487 & 3.51$\times$\\}
\newcommand{\HeadlineConclusion}{The degree-12 three-bin prototype improves the accepted DSpec calculation by 3.05 times for one objective and 4.43 times for two. Its measured times are 0.240/0.316 s versus YSpec's 0.266/0.312 s. The smaller-degree case remains slower than YSpec. These are measured single-thread prototypes, not changes to the production default.}
\newcommand{\SpectralRows}{3 & 3 & 0.42 & 0.67 & 0.29 & 5.15 & 1.70\\
3 & 48 & 6.56 & 0.77 & 0.81 & 4.96 & 1.75\\
6 & 3 & 4.00 & 11.81 & 1.78 & 121.15 & 23.41\\
6 & 48 & 66.31 & 14.41 & 5.60 & 149.94 & 30.07\\
12 & 3 & 9.60 & 50.72 & 19.69 & 1067.08 & 100.53\\
12 & 48 & 119.74 & 52.96 & 51.05 & 962.89 & 106.33\\}
\newcommand{\SpectralPlanTwelve}{0.96}
\newcommand{\AblationText}{\begin{center}
\begin{tabular}{lrl}
\toprule
Actual L12/F3/full-grid route & Median (s) & Frozen volume gates\\
\midrule
v4 real & 0.313 & pass\\
v5 real & 0.289 & pass\\
v5 complex & 0.231 & pass\\
packed scalar & 0.457 & pass\\
packed eigen & 0.529 & pass\\
packed openblas & 0.513 & pass\\
minimal basis & 0.275 & pass\\
rstrain bridge & 0.337 & pass\\
strict native build & 0.275 & fail\\
\bottomrule
\end{tabular}
\end{center}
The complex tensor layout is faster here than the real-degree split, despite its higher nominal arithmetic count. Small packed Eigen/OpenBLAS calls lose to the packed scalar route; matrix-matrix work over many radii pays. The nodal r-strain bridge loses on this grid. Strict native-architecture compilation (without fast-math, FMA contraction disabled) fails the processed-spectrum ($1.42\times10^{-12}$) and objective ($6.46\times10^{-12}$) gates and is rejected. The final headline table therefore uses portable O3 and the complex tensor layout.

The v6 forward angular cache adds actual reuse across objectives: it retains the volume forward strains and receiver displacement by exact grid/frequency identity. Hits skip source tensor synthesis and receiver suffix sampling; workspace preparation and representative coefficient gathering still run. The separate radial response cache saves repeated solves. At L12/F3 the forward field payload is 21.8 MB. Zero-budget fallback and contact-bearing reuse with radial response caching disabled both pass frozen gates. This scheduling/cache prototype is distinct from a general multi-objective API.
}
\newcommand{\ValidationText}{The final v6 headline outputs pass all 18 independent process/configuration comparisons against the unchanged production evaluator: returned raw/processed spectra, objectives and all six stored modulus channels, at exact matched coordinates. Repeated outputs are deterministic. Maximum kernel relative L2 is $8.12\times10^{-14}$ and scaled pointwise $9.04\times10^{-12}$. Expanded tests separately establish 336 public field and 240 source-contact comparisons per bridge setting against the saved production arrays. An additional complete permuted PREM tensor product agrees with the standard field reconstruction in solid strains/fluid divergence below $1.13\times10^{-15}$; it is not the source of those historical contact arrays.

Fresh v5 optimized tensor Born checks pass the independent Cartesian regular and source-local volume/contact/total gates, with the existing factorization and radial-response cache disabled. The v6 forward angular cache is checked separately against frozen full outputs; it was not used in that Born adapter run. The final forward-cache tests pass all eight frozen comparisons at default and zero budgets, covering two/three objectives and complex/real layouts. Specifically, contact-bearing L3/three-objective and L12/two-objective cases use the complex layout and radial-response cache off; L12/three-objective complex/real cases use radial-response cache on and contacts off. The response-cache contacts-on and native L12 failures remain explicitly rejected; all gates were retained.
}
\newcommand{\MemoryText}{For degree 12, one/two objectives, maximum sampled process high-water memory is Phase 1: 44.5/45.0 MiB; prototype: 83.9/124.2 MiB; YSpec: 17.9/27.7 MiB. These values are polled after executable entry every 10 ms and include setup, outputs and, in the prototype driver, its subsequent same-engine reference calculation. They are process-level diagnostics, not isolated API allocator peaks or aggregate limits. The runtime advantage has a clear memory cost. The tensor coefficient multiplication is tiled, but full volume strain/output arrays and plan maps are still materialized; production integration should bound those intermediates more tightly.
}

\textbf{Finding.} The remaining cost is principally an angular-synthesis and data-reuse problem. It is not explained by interpolation or by the final six strain products. The implemented tensor prototype reuses angular plans and evaluates many radii with matrix products. The unchanged production implementation and YSpec are compared below. The matched model has full gravity, no attenuation, NN=8 SEM interpolation, a 50 km-deep source at $10^\circ/20^\circ$, and a 3 km-deep receiver at $45^\circ/90^\circ$. All new implementations are isolated diagnostic prototypes; production adoption is a separate decision.

\section{Matched end-to-end comparison}
\begin{center}
\begin{tabular}{rrrrrrr}
\toprule
$L$ & Bins & Objectives & Phase 1 (s) & Prototype (s) & YSpec (s) & Gain\\
\midrule
\HeadlineRows
\bottomrule
\end{tabular}
\end{center}
Each entry is the median of three fresh processes, alternating solver order, with one numerical thread. The volume grid contains 56 element/radius traces and 192 angles at $L=3$, or 1,352 angles at $L=6,12$: respectively 10,752 or 75,712 points. API times exclude model/SEM setup and writing output, but include first-use plans and measurement processing. The machine is an Intel Core Ultra 5 235U; native headline runs share one fixed allowed CPU. Both codebases use portable GCC O3/Release without fast-math. Two DSpec objective calls are summed; YSpec handles both in one batch.\footnote{The prototype's radial response and forward angular caches are enabled for the two-objective, volume-only case. Source contacts are excluded from every timing in this table. Contacts are independently checked with radial-response caching disabled. Timing identities, commands, ranges, sampled process memory and output arrays are saved in \texttt{final-timings-v6/}.}

The three-bin benchmark solves through \textbf{1.171875 mHz}; 12 and 48 bins extend to 4.6875 and 18.75 mHz. The original nominal 0.1--0.8 mHz measurement band is not the maximum solved raw frequency. Extended cases are fixed-discretization performance/equivalence probes; they establish no high-frequency or angular convergence. YSpec requires at least two bins, so the one-bin case is DSpec-only.

\textbf{Practical result.} \HeadlineConclusion

\section{Why ``two solves and multiplications'' was incomplete}
The adjoint formulation needs a forward field and an adjoint field for each independent objective. In this implementation, obtaining either field means solving one radial system for each degree, family and frequency. At $L=12$, one frequency has 12 spheroidal, 12 toroidal and one radial systems. Three bins therefore require 75 distinct dynamic matrices. The accepted Phase 1 route factors each matrix twice, giving 150 factorizations; the prototype shares one factorization between forward and Hermitian-adjoint substitutions. A changing frequency changes the matrix and requires new numerical factors.

There are 75,712 spatial points and six strain slots per field, so one field/frequency has 454,272 complex strain outputs. Reconstructing those outputs also requires summing degree/order coefficients. The profiled calculation performs about \textbf{128 million complex strain multiply-accumulates}, across two fields and three bins. The final pointwise products are cheap; the earlier explanation conflated that cheap final work with the harmonic reconstruction.

\begin{center}
\begin{tabular}{p{.56\textwidth}rr}
\toprule
Accepted Phase 1, $L=12$, 3 bins, 1 objective & Seconds & Calls/misses\\
\midrule
Degree sums and order phases & 0.2683 & 6,702\\
Wigner table construction & 0.1455 & 1,989\\
Geometry grouping & 0.0760 & 120\\
Factorization & 0.0779 & 150\\
Angular coefficient table construction & 0.0377 & 1,971\\
Global weighted-adjoint substitution & 0.0230 & 72\\
System matrix formation & 0.0185 & 150\\
Rotation/addition scratch & 0.0146 & 2,229\\
Radial coefficient preparation & 0.0062 & 342\\
Final kernel products & 0.0026 & 57\\
\bottomrule
\end{tabular}
\end{center}
These are instrumented scopes and some overlap: do not sum them into a total. The native accepted API took 0.7004 s; instrumentation took 0.7423 s. Output arrays were byte-identical. Table construction was repeatedly missing the cache because a 4,096-point block split the shared angular set differently on successive radial traces. This is an avoidable algorithm/layout defect, rather than an inherent cost of adjoints.

\section{What YSpec actually does differently}
Inspection of \texttt{SampledKernels.cpp} confirms that YSpec prepares strain coefficients for all sampled radial nodes, with degree contiguous and the real/imaginary degree dot products explicit. It groups 64 angular places in a block, constructs source/receiver Legendre tables for that block, and reuses them over \emph{all radial nodes}. Geometry is prepared before the frequency loop. The inspected version rebuilds those block Legendre tables at each frequency; it does not already implement every cross-frequency reuse proposed here.

YSpec constructs a basis of measurement directions and shares its radial responses and field products across objectives. It contracts in the source frame, applying the relative source/receiver spin rotation to the adjoint. Its radial boundary-value solver and DSpec's assembled SEM solve are different methods, so equal physics and output samples do not imply equal linear algebra or degrees of freedom. The informative weighted-profile/integral comparison at the existing $10^{-3}$ tolerance passes for the established small and large benchmark objectives. It does not certify whole-planet receiver-region strain accuracy.

\begin{center}
\begin{tabular}{p{.25\textwidth}p{.32\textwidth}p{.33\textwidth}}
\toprule
Work & Accepted DSpec Phase 1 & Implemented audit alternative\\
\midrule
Angular addressing & Spatial chunks can fragment one angle set & Exact trace $\times$ angle plan, preserving output order\\
Angular synthesis & Contiguous degree dots per radial trace & Real or complex matrix products over 16 radii\\
Angular plans & Chunk-keyed reuse often misses & Persistent frame/grid/spin/degree tables\\
Radial factors & Separate forward/adjoint factorization & Shared equilibrated factorization\\
Multiple objectives & Repeated scalar evaluations & Cached source and minimal receiver responses\\
Output radii & General element-side radii & General or native GLL traces, separately tested\\
Frequency order & Synthesize both fields for every bin & Exact same-frequency coefficient covariance\\
\bottomrule
\end{tabular}
\end{center}

\section{Implemented direct routes and their costs}
The prototypes are real executable routes consuming the SEM solutions. They include packed scalar, Eigen and one-thread OpenBLAS synthesis; complex tensor products; real degree-wise tensor products followed by each order phase; shared factors; minimal supplemental receiver basis construction (weighted global reconstruction is used only in response-cache mode); repeated-objective response reuse; direct arbitrary-radius interpolation; native differentiation; and an optional nodal strain-numerator bridge. Diagnostic packet export is disabled during timings.

The final tensor plan stores the exact point-address map, transformed angles, spin rotations, Wigner values and real/complex synthesis bases. Its key distinguishes frame, grid ordering, degree and coefficient labels. It persists across frequency and objective loops, rather than across mutable model updates. A complete permuted tensor grid is accepted and restored to input order; duplicate or ragged angle sets are rejected by this explicit route. Packed sampling remains available for arbitrary point clouds.\footnote{The first tensor prototype rebuilt its main angular tables every call despite caching representative coefficient extraction. Review caught this; its v4 binary and evidence remain saved, and v5 adds actual reuse.}

Only the six strain slots are retained on tensor volume points; full field records are still used for a small representative set and receiver displacement suffix. Fluid coefficient synthesis uses divergence, and forbidden $l<|s|$ spin slots vanish. Both completed fields are still rotated for the direct point-kernel route. A separate relative-spin helper is not an independently integrated direct backend; relative angular factors are included in the spectral product representation. The measured baseline rotation cost is much smaller than degree summation and table construction.

\begin{center}
\begin{tabular}{p{.56\textwidth}r}
\toprule
Native accepted-engine block/cache control & Median API (s)\\
\midrule
4,096-point blocks, 32 MiB payload budget & 0.7139\\
1,352-point blocks, 32 MiB & 0.5497\\
2,704-point blocks, 32 MiB & 0.5343\\
One 75,712-point block, 32 MiB & 0.5792\\
4,096-point blocks, 128 MiB & 0.6451\\
\bottomrule
\end{tabular}
\end{center}
All five controls produced byte-identical outputs, each in three processes. The result shows that matching chunk boundaries to semantic angular structure helps; merely enlarging caches or materializing everything is less effective. These controls are separate from the final isolated comparison and should not be combined as multiplicative speedups.

\AblationText

\section{Native SEM output and interpolation}
There is no mathematical requirement to request arbitrary output radii. Native one-sided GLL traces are a natural kernel representation. At a node, values are a gather and derivatives follow from $D_e=D_{\mathrm{ref}}/J_e$. Distinct element sides must remain distinct even when their radii coincide; source-lift ownership and fluid/solid traces must not be averaged.

For arbitrary radii, storing a derivative polynomial does not make evaluating it at a new coordinate free. The interpolation weights can be prepared once per exact element/radius and applied to all relevant solution columns. More importantly, radial coefficient preparation is only about 6 ms of the profiled 742 ms. Removing it alone cannot explain or close the main performance gap.

The implemented bridge forms the six nodal $r\varepsilon$ numerators from $U,V,W$ and their derivatives. For element degree $p$, each numerator has degree at most $p$, so interpolating it and then dividing by $r>0$ is algebraically exact. Interpolating raw nodal strain directly would lose its $1/r$ dependence. Forming all nodal numerators can cost more than direct evaluation for sparse output radii; it is an optional representation with a density-dependent crossover.

\textbf{Numerical boundary.} Native $L=3$ passed the new frozen comparison, but the $L=12$ all-element GLL output exceeded the unchanged $10^{-12}$ relative-L2 kernel gate: about $2.69\times10^{-12}$ with Dref/J, and $1.69\times10^{-12}$ with cached physical-coordinate nodal derivatives. The scaled pointwise gate passed. This is a small discrepancy, but the branch is not accepted under the current contract. The physical derivative control reduced it without establishing a complete cause. General-radius preservation passes; replacing the default output grid is deferred until this discrepancy is resolved.

A further dense-output representation is theoretically possible: $r^2K$ has degree at most $2p$ inside an element, so accumulate its polynomial before sampling many radii. It requires at least $2p+1$ suitable values or coefficient products, not merely $p+1$ native kernel values. The present grid has roughly eight traces per element versus $2p+1=15$ at NN=8, so this is not a promising first implementation for the measured workload. It remains relevant to much denser requested radial output.

\section{Frequency-first contraction: the algorithmic change}
Let $f_k$ and $a_k$ be compact source- and receiver-frame strain coefficient vectors at one radial trace and frequency $k$, and $h_F(x),h_A(x)$ their angle-dependent basis vectors. The exact rearrangement is
\begin{align*}
 C&=\sum_k\overline{a_k}\,f_k^T,\\
 \sum_k\overline{a_k(x)}f_k(x)
 &=\overline{h_A(x)}^T C h_F(x).
\end{align*}
Every same-frequency $l,l',m,m'$ pair is retained. Multiplying separately frequency-summed fields would introduce cross-frequency terms and be wrong. Summing kernels per degree would discard cross-degree terms and also be wrong. The actual measurement adjoint $P^*q$ remains a full-frequency operation including the Hann window, damping and SI factors.

At $L=12$, the local source and receiver widths are approximately $q_F=59$ and $q_A=37$, including $l=0$. Direct synthesis costs $O(FP(q_F+q_A))$; covariance formation/evaluation costs $O(Fq_Fq_A+Pq_Fq_A)$ per pair and radius. The arithmetic frequency crossover is around 23 bins. This explains why the three-bin case favors direct synthesis and why the 48-bin case is worth implementing. Actual constants and the quality of the direct matrix implementation matter.

The standalone consumer uses actual solver-derived local coefficients, including spheroidal/toroidal sums, normalized radial $l=0$, and exact source/receiver frame information. All nine $L=3,6,12$ and $F=3,12,48$ cases were evaluated over \textbf{three complete radial traces}, with 192 or 1,352 angles each. It reconstructs all six strain slots and checks the resulting modulus products against direct pointwise synthesis. Maximum kernel relative L2 is $3.42\times10^{-15}$; maximum scaled pointwise difference is $2.46\times10^{-13}$. These are volume equivalence checks, not contact or whole-engine timing acceptance.

\begin{center}
\begin{tabular}{rrrrrrr}
\toprule
$L$ & $F$ & Direct (ms) & Covariance (ms) & Recouple/FFT (ms) & Plan (ms) & SVD (ms)\\
\midrule
\SpectralRows
\bottomrule
\end{tabular}
\end{center}
Entries are medians of three independent processes on exactly three traces. Direct means per-frequency coefficient synthesis; covariance includes formation and angular evaluation. Recouple/FFT includes common-frame coefficient rotation/covariance recoupling, scalar kernel assembly and inverse-FFT synthesis. The full combined validator plan is shown separately. At L12/F48 its compact angular basis costs about 66 ms, global frame/QR about 822 ms, Gaunt construction about 66 ms, and grid addressing 0.1 ms. Covariance needs the compact basis, not the global/QR/Gaunt stages; a cold covariance request must be charged that appropriate compact setup. FFTW plan creation is inside the FFT projection/synthesis timers. Radial solves, measurement processing, coefficient export/parsing and contacts are outside these columns. No three-trace result is extrapolated into a measured full-volume API time. A separate v5 real-tensor run on the matching three traces takes about 20 ms sampling at F3 and 245 ms at F48, but it includes coefficient extraction and receiver work. The consumer's direct primitive is not the final optimized complex tensor backend, so its covariance speedup must not be advertised as an end-to-end gain over that backend.

\section{Low rank, common-frame harmonics and FFT}
Each complex covariance has rank at most $\min(F,q_A,q_F)$ and is generally neither Hermitian nor positive semidefinite. The implemented thin SVD verifies candidate ranks against actual output L2 and pointwise gates instead of choosing a singular-value cutoff. Observed ranks at 48 bins are 2--4 ($L=3$), 5--7 ($L=6$), and 11--13 ($L=12$). Lower factor storage does not guarantee lower total cost: SVD and verification lose to the direct covariance route on the measured workloads.

Products of degree-$L$ fields have scalar harmonic bandlimit $2L$. A common-frame representation therefore needs at most $(2L+1)^2=625$ coefficients at $L=12$; the source-frame azimuthal restriction can reduce that to 515. Both compact pole frames cannot simultaneously retain their small order sets after rotation to a common frame. The prototype performs verified coefficient-frame conversion, quadrature-defined Gaunt recoupling through $2L$, and global longitude FFT synthesis. This avoids guessing phase or normalization conventions.\footnote{Primary harmonic product/Gaunt identities: \href{https://dlmf.nist.gov/34.3}{NIST DLMF 34.3}. The pinned GSHTrans convention is orthonormal. A prior normalization assumption and two Eigen temporary-expression lifetime bugs were found and fixed; failed outputs remain in the evidence log.}

The $26\times52$ global angle grid resolves degree-24 kernel output. Gaunt and FFT coefficient comparisons have maximum relative L2 $1.94\times10^{-14}$ and scaled pointwise $2.06\times10^{-12}$; FFT and direct harmonic synthesis agree below $1.51\times10^{-13}$ scaled pointwise. Sparse points use covariance/direct evaluation; FFT requires a suitable complete tensor grid.

The full recoupling plan uses numerical frame conversion and quadrature Gaunt construction, and costs about \SpectralPlanTwelve{} s at $L=12$. An analytic Wigner rotation/Gaunt implementation or saved validated plan could reduce that cost; no speedup for such an implementation has been measured. The present warm 48-bin recoupling/FFT path is comparable to covariance at $L=12$. Its rough explicit payload is 154.7 MB at L12: 12.3 MB compact bases, 66.7 MB global/frame data, 75.5 MB Gaunt data, and 0.14 MB grid addressing, excluding allocator overhead. It is a viable exact representation, not the current default recommendation.

\section{Radial solver alternatives: implemented and rejected}
The SEM operator has narrow bandwidth in physical indexing (ordinary spheroidal half-bandwidth at most about 23 for NN=8). Implemented alternatives include equilibrated pivoted general-band LU, NaturalOrdering SparseLU with verified symbolic-pattern reuse, and element-interior static condensation with full interior recovery. Forward and Hermitian-adjoint substitutions use the same factors. Condensation uses actual LTG/interface identity, pivoted local blocks and Schur solves. It checks local invertibility; global solvability alone does not guarantee that eliminated blocks are safe.\footnote{Band packing and conjugate-transpose solve contracts follow the primary Netlib \href{https://www.netlib.org/lapack/complex16/zgbtrf.f}{ZGBTRF} and \href{https://www.netlib.org/lapack/complex16/zgbtrs.f}{ZGBTRS} sources.}

The full PREM $L=12$, bins 1--48 run contains 4,800 backend records. Original-system residuals are small (forward around $5\times10^{-17}$, adjoint around $2\times10^{-20}$), but this does \emph{not} establish output equivalence. The changed elimination orders fail the frozen field/contact gates:
\begin{center}
\begin{tabular}{lrrr}
\toprule
Backend & Failed field groups & Failed contacts & Max contact relative L2\\
\midrule
Band LU & 252 & 106 & $5.38\times10^{-9}$\\
NaturalOrdering & 246 & 106 & $6.78\times10^{-9}$\\
Static condensation & 258 & 104 & $8.36\times10^{-10}$\\
\bottomrule
\end{tabular}
\end{center}
Counts are from one bridge setting, 336 field and 240 contact groups each; both bridge settings retain the same disposition. Small residuals and overall solution-vector norms do not bound derivative-sensitive field/contact cancellations. None of these factor replacements is accepted. The shared existing equilibrated solver passes the expanded frozen field/contact comparison and is the supported optimization candidate.

\section{Acceptance, memory and provenance}
\ValidationText

All thresholds are unchanged: $10^{-12}$ field norms except $10^{-9}$ adjoint strain; $10^{-12}$ kernel L2 and $10^{-10}$ scaled pointwise with the existing 1\%-of-peak floor; $10^{-12}$ contacts; informative YSpec weighted profiles/integrals $10^{-3}$. Tiny physical SI kernels are compared against their actual norms, with no unit-sized denominator that could hide errors.

The plan's 32 MiB and response cache's 512 MiB limits are \emph{retained-payload admission budgets}. They omit map/key/allocator overhead, temporary matrices, deep-copy return values and outputs. Capacity or response-span failures in the current diagnostic response cache throw; there is no automatic direct-solve recovery. Its exact frequency/RHS keys are valid for an immutable live SEM object, and it must not survive operator mutation. The batch option is explicit opt-in: enabling contacts does not automatically bypass it. Contacts-on batch results fail some frozen checks, while tested volume-only results pass.

\MemoryText

Protected production files and existing Phase 0 evidence are verified against the pre-audit manifest. YSpec remains clean, both repository HEADs are unchanged, and no commit or push was made. Compiler concurrency was limited to two jobs; numerical calculations used one thread. This is bounded runtime control and sampled process diagnostics, not kernel aggregate-memory accounting. The existing 177-test Phase 1 regression evidence remains distinct from the fresh prototype field/kernel/Born checks; the full production suite was not rerun for isolated prototypes. Spectral v2 is the reproducible current-API source/binary. The older v1 snapshot contains an incompatible Wigner template alias and is superseded; its historical timings are not used here. Fresh v2 solver-packet oracles and ranks pass after the API correction.

\section{Disposition and concrete adoption sequence}
\begin{enumerate}
\item \textbf{Promote semantic angular planning and tiled direct synthesis first.} Keep the existing equilibrated factorization and arbitrary-radius API. Integrate one shared factor per matrix, exact receiver basis handling, persistent grid/frame plans and the measured best real/complex matrix layout. Preserve the packed fallback and strict contact path. Repeat the existing production regression gates on the integrated change.
\item \textbf{Integrate objective reuse with separate contact rules.} Preserve each full $P^*$ action. Forward angular-field reuse has passed contact-bearing checks with the original weighted solves. The batch response cache saves repeated radial solves, but must gain an explicit contact-safe route and robust bounded-cache fallback before becoming a general API default.
\item \textbf{Offer exact frequency covariance for sufficiently many bins.} Integrate the proven covariance algorithm behind a measured direct/covariance dispatch; measure the complete solver-to-output path and tile coefficient history by radial blocks. The standalone warm advantage is established; a full-grid end-to-end gain is not yet measured.
\item \textbf{Retain native SEM output as a deliberate feature.} Resolve the strict $L=12$ native-node discrepancy before making it the default. Retain optional arbitrary-radius resampling and side provenance. Do not spend the main performance effort on interpolation alone.
\end{enumerate}

No approximate degree decoupling, cross-frequency product, unverified rank truncation, Hermitian-positive-definite solver, raw-strain interpolation, or negative-order real-field symmetry is a valid shortcut here. GPU and multi-thread designs need a separately authorized hardware/resource comparison; this audit supplies the single-thread algorithmic baseline. Strict-IEEE architecture-specific compiler tuning is a useful secondary control, rather than a substitute for the demonstrated data reuse. No fast-math flags were used.

\textbf{Coverage boundary.} Promising mathematical families were derived and runnable prototypes were exercised before disposition. This is not a claim that every possible implementation or hardware backend has been exhausted. Receiver-region physical limitations, mesh/degree convergence, general attenuation/measurement combinations and large receiver ensembles require their own validation. Production numerical code remains the accepted Phase 1 version; the audit delivers concrete implementations, preserved failures and a measured integration order.

\clearpage
\section{Independent review and evidence index}
Fresh independent Luna High review and subsequent fresh Sol High review both pass the bounded prototype audit. They independently checked the final source/binary identities, saved numerical gates and benchmark summaries; neither reviewer changed numerical sources or ran the numerical campaign. Their verdicts preserve the rejected paths and require staged production integration with regression checks.

The evidence directory is \texttt{investigations/dspec-sk-profile-20261026/algorithm-audit/} in this chat workspace. Relative paths below locate the runnable implementations and authoritative evidence; preliminary and rejected attempts are retained separately.
\begin{center}
\small
\begin{tabular}{p{.39\textwidth}p{.53\textwidth}}
\toprule
Evidence & Path\\
\midrule
Scope, resources and frozen gates & \texttt{protocol.json}\\
Original protection/final verification & \texttt{before.json}, \texttt{final\_protections.json}\\
Theory before prototypes & \texttt{theory-before-implementation.md}, solver addendum\\
Accepted implementation profile & \texttt{current\_profile\_summary.json}\\
Final direct candidate/compile hashes & \texttt{fast/v6/}, \texttt{freeze-v6.sha256}\\
Matched final API times/output gates & \texttt{final-timings-v6/}\\
Layout/radial/compiler controls & \texttt{final-ablation-v5/}\\
Frozen candidate/path-specific failures & \texttt{final-validation-v6/}\\
Forward reuse and zero-budget checks & \texttt{forward-cache-gates-v6/}\\
PREM tensor/field/contact checks & \texttt{field-validation-v5-tensor/}\\
Fresh independent tensor Born oracle & \texttt{born-v5/}\\
Alternative factors and residuals & \texttt{solvers/evidence\_manifest.json}\\
Reproducible spectral consumer & \texttt{spectral/frozen/v2/}\\
Spectral repeated costs & \texttt{spectral/isolated-repetitions-v2/}\\
Spectral 12-bin controls & \texttt{spectral/v2-f12-checks/}\\
Independent review records & \texttt{reviews/}\\
\bottomrule
\end{tabular}
\end{center}
The final headline driver syntax is \texttt{fast\_driver-v6 MODEL SAMPLES OUT L F NOBJ BACKEND RADIAL SOLVER}; the output directory must be empty. Exact commands and environment switches for every accepted/rejected case are recorded in the run manifests. Reproduction should use a fresh evidence destination and the frozen source build command, preserving the existing production checkout and original evidence.
